%!TEX root = ../main.tex
% =============================================================================
%  八、结果与分析                建议 4–5 页
%  评审在问：数对不对？有没有解释？
%  本节安排：① 口径框 C-02、C-03、C-04、C-06、C-07、C-08（编号按交底书 §6.1 唯一分配；
%               C-01、C-05、C-09、C-10、C-11 分别首次出现于第四、五、六、七章）；
%            ② 四问结果总表（含无储能基线与节省率）；
%            ③ 成本分解（fig05 为问题 1 日费用瀑布）与策略机制解释（fig03、fig10）；
%            ④ 题面表 1、表 2 格式：问题 2（主模型）在四个指定日期上完整给出；
%            ⑤ 题面表 3 格式：四种方案 × 四个指定日期合并为一张规整表；
%            ⑥ 问题 3、4-2、4-3 的四个指定日期汇总表 + 结果文件指引（附录 C 有对照表）。
%  红线：本节每一个数值都必须与 result*.xlsx 逐格一致（见第九章审计 A9、A10）。
% =============================================================================

\section{结果与分析}

本节给出四问的数值结果与解释；数字取自 \texttt{intermediate/paper\_numbers.json}
（\texttt{code/paper\_numbers.py} 汇总），口径由 \texttt{code/config.py} 唯一确定，
可由 \texttt{code/audit.py} 独立复算（见第九章）。

\subsection{指标口径表}

本节把在本题参数范围内需要逐项说明的六个核心指标口径一次列清，其余引用其定义章：
无储能基线 \calref{C-01}（第四章，V1 固定价窗口／V2 波动价窗口／V3 典型日）、
紧急购电量 \calref{C-05}（第五章）、调整与违约费用 \calref{C-09}、信息价值 \calref{C-10}（第六章）、
条件风险价值 \calref{C-11}（第七章）。

\begin{caliber}{C-02}{全天购电量}
\textbf{定义}\quad \(E^{\text{day}}=\sum_{t=1}^{144}E_t\)，即该日 144 个 10~min 槽的\textbf{计划购电量}之和；%
\textbf{公式}\quad 结果为逐槽求和，无闭式；结果文件每行末列即该日汇总；%
\textbf{数据}\quad 结果文件「计划购电量」工作表的 144 个时段列；%
\textbf{时段}\quad 单日；问题 1 为典型日，问题 2--4 为 2025-02-01 至 12-31 的 334 天；%
\textbf{单位}\quad kWh（结果文件保留 4 位小数，正文保留 2--4 位）；%
\textbf{含}\quad 问题 2、4-2 中即第一阶段计划量；问题 3、4-3 中调整后的量一律称"全天调整购电量"并单独命名；%
\textbf{不含}\quad 不含储能充电量、不含紧急购电量 \calref{C-05}、不含盈余处置量 \calref{C-06}；%
\textbf{对照}\quad 问题 2、3 的计划购电量模型内值为 \(22\,220\,333.2033\)~kWh，结果文件回读值为 \(22\,220\,333.1391\)~kWh，
相对差 \(2.9\times10^{-9}\)，来自写入时的 4 位小数取整；%
\textbf{复算}\quad \texttt{results\_io.py::plan\_sheet\_frame()}（写出行汇总列）与 \texttt{audit.py::audit\_result\_files()}（回读结果文件后重算 \(\sum_{144}\) 槽）
\end{caliber}

\begin{caliber}{C-03}{总购电费用（计划 \(+\) 紧急 \(+\) 调整相关）}
\textbf{定义}\quad \(C^{\text{total}}=C^{\text{plan}}+C^{\text{breach}}+C^{\text{excess}}+C^{\text{emg}}\)；
问题 2、4-2 无调整机制，后两项恒为 0；%
\textbf{公式}\quad 问题 2、4-2：\(C=\sum_tp_tE_t+5\sum_tp_tg_t\)；
问题 3、4-3：\(C^{\text{plan}},C^{\text{breach}},C^{\text{excess}}\) 按 \calref{C-09} 的偏差量定价口径计算，
\(C^{\text{emg}}=\sum_t5p_tg_t\)；%
\textbf{数据}\quad 电价取附件 1（问题 1--3）或附件 4（问题 4）；电量取本队结果文件；%
\textbf{时段}\quad 问题 1 为典型日；问题 2--4 为 334 天；%
\textbf{单位}\quad 元（保留 2 位小数，内部计算不取整）；%
\textbf{不含}\quad 不含储能投资与运维成本、容量补偿、辅助服务收益、绿电环境价值、循环衰减与安全成本；%
\textbf{对照}\quad 若只报计划购电费而不计紧急购电，问题 2 为 \(13\,637\,552.90\)~元，
比 \(C^{\text{total}}\) 低 \(627\,239.24\)~元（\(4.40\%\)）；%
\textbf{复算}\quad \texttt{q2.py::run\_q2()}、\texttt{q3.py::summary()}、\texttt{q4.py::run\_q4\_2()} 与 \texttt{audit.py::audit\_result\_files()}，
四项之和与总费用的偏差为 0
\end{caliber}

\begin{caliber}{C-04}{节省率}
\textbf{定义}\quad \(\eta=\bigl(C^{base}-C^{\text{total}}\bigr)/C^{base}\times100\%\)，分母与分子必须同窗口、同电价口径；%
\textbf{公式}\quad \(C^{base}\) 取 \calref{C-01} 的对应变体（V1 固定价窗口 / V2 波动价窗口 / V3 典型日）；%
\textbf{数据}\quad \(C^{base}\) 与 \(C^{\text{total}}\) 均取自本队结果；%
\textbf{时段}\quad 与 \calref{C-01} 变体一致；%
\textbf{单位}\quad \%（保留 2 位小数）；%
\textbf{不含}\quad 不得相对贪心策略或完美信息下界报同一个数而不加限定；相对其他基线的结果必须\textbf{单独命名}；%
\textbf{对照}\quad 问题 1 相对 V3 为 \(26.90\%\)；问题 2、3 相对 V1 为 \(13.06\%\)、\(15.08\%\)；
问题 4-2、4-3 相对 V2 为 \(11.35\%\)、\(13.54\%\)；%
\textbf{复算}\quad \texttt{paper\_numbers.py::main()}（由 \calref{C-01} 与 \calref{C-03} 导出）
\end{caliber}

\begin{caliber}{C-06}{弃光/盈余处置电量}
\textbf{定义}\quad 光伏盈余与超购部分的处置量：典型日按
\(s_t=PV_t\Delta t+E_t+d_t-L_t\Delta t-c_t\ge0\) 逐槽统计（式~\eqref{eq:q1-balance} 的松弛量）；
窗口口径按实际数据复算为 \(\sum_t\max\{PV_t\Delta t+E_t+d_t+g_t-L_t\Delta t-c_t,0\}\)；%
\textbf{公式}\quad 上式为定义式；该量非负，其余额为超购部分被弃置的电量；%
\textbf{数据}\quad \(L,PV\) 取附件 2 实际值，其余取本队结果文件；%
\textbf{时段}\quad 逐槽统计、按窗口（334 天）汇总；%
\textbf{单位}\quad kWh（按电量统计，\textbf{不得}按功率报）；%
\textbf{不含}\quad 不含"因充电功率受限而未发生的充电"，不含未购电部分；%
\textbf{对照}\quad 问题 1 典型日最优弃光 \(=0.0000\)~kWh（数值残差 \(1.37\times10^{-12}\)~kWh）；
问题 2、3 的窗口盈余处置量为 \(3\,106\,390.84\)~kWh 与 \(2\,638\,963.35\)~kWh；
问题 4-2、4-3 的同一量尚未纳入汇总；%
\textbf{复算}\quad \texttt{config.py::curtailment(price, load, pv, E, c, d)}（典型日）与 \texttt{audit.py::audit\_internal()}（窗口口径）
\end{caliber}

\begin{caliber}{C-07}{充放电量与储电量（交流侧）}
\textbf{定义}\quad 结果文件中的 \(c_t,d_t\) 为\textbf{并网点（AC 侧）}电量；储电量为\textbf{时段末}值 \(S_t\)；%
\textbf{公式}\quad 递推见式~\eqref{eq:q1-soc}，\(\eta_c=\eta_d=0.9\)（单边口径，往返 \(0.81\)）；
充放功率上限 \(5000\)~kW 对应单槽上限 \(\overline P\Delta t=833.3333\)~kWh；%
\textbf{数据}\quad 本队结果文件「充放电量」工作表；%
\textbf{时段}\quad 逐槽；表 2 格式按 6 个 4~h 块（每块 24 槽）汇总，并给出 \(0{:}00\) 与 \(24{:}00\) 两个储电量点；%
\textbf{单位}\quad kWh；运行区间 \([1200,10800]\)~kWh；%
\textbf{不含}\quad 不含电池侧（DC）电量、不含自放电与循环衰减（假设 A4）；%
\textbf{对照}\quad 恒等式校验：问题 1 的充放比 \(20\,740.67/16\,799.94=1.2346=1/(\eta_c\eta_d)\)，
与 \(\eta_c\eta_d=0.81\) 互为倒数，偏差 \(4.44\times10^{-16}\)；问题 2 四个指定日期的合计充放比同为 \(1.2346\)；%
\textbf{复算}\quad \texttt{results\_io.py::soc\_sheet\_frame()}（块级汇总）与 \texttt{audit.py::audit\_result\_files()}（块级净递推）
\end{caliber}

\begin{caliber}{C-08}{储电量递推与终端条件}
\textbf{定义}\quad 储电量\textbf{跨日连续递推}（\(S_{d,0}=S_{d-1,144}\)，假设 A5），
窗口末端施加终端条件 \(S_T\ge S_0\)（假设 A6），问题 1 另加严格日周期 \(S_0=S_{144}\)；%
\textbf{公式}\quad \(S_t=S_{t-1}+\eta_cc_t-d_t/\eta_d\)，\(S_t\in[1200,10800]\)~kWh；
2025 年 1 月为预热期（假设 A8），其终值即 2 月 1 日初值；%
\textbf{数据}\quad 结果文件「充放电量」工作表的储电量列与中间结果 \texttt{q2\_main.json} 等；%
\textbf{时段}\quad 逐槽；窗口 2025-02-01 至 12-31（334 天）；%
\textbf{单位}\quad kWh；%
\textbf{不含}\quad 不含"每日清零""每日重置为 6000"两种读法——前者会引入末日掏空型短视解，后者与 1 月预热期矛盾；%
\textbf{对照}\quad 问题 1：日周期口径 \(35\,126.9486\)~元（\(S_0=S_{144}=6000.00\)~kWh）；
仅要求两端相等的自由环状口径 \(35\,118.5986\)~元（最优 \(S=8550.00\)~kWh），相差 \(8.35\)~元（\(0.02\%\)）；
窗口口径下"末值低于初值的天数"为 0；%
\textbf{复算}\quad \texttt{lp\_engine.py::solve\_day()}（终端约束）与 \texttt{audit.py::audit\_result\_files()}（\texttt{A6\_末值低于初值天数}）
\end{caliber}

\subsection{四问结果的汇总与对照}

\begin{table}[htbp]
  \centering
  \caption{四问主要结果汇总（费用单位：元，电量单位：kWh；问题 1 为典型日，
           其余为 2025-02-01 至 2025-12-31 共 334 天；问题 4 的电价为附件 4 波动价；
           表中带 a 标记者见表注）}
  \label{tab:summary}
  \scriptsize
  \setlength{\tabcolsep}{3pt}
  \begin{tabularx}{\linewidth}{L{2.35cm} C{1.60cm} C{1.95cm} C{1.95cm} C{1.95cm} Y}
    \toprule
    指标 & 问题 1 & 问题 2 & 问题 3 & 问题 4-2 & 问题 4-3 \\
    \midrule
    计划购电量 & 59\,482.70 & 22\,220\,333.14 & 22\,220\,333.14 & 22\,198\,236.57 & 22\,255\,478.92 \\
    调整购电量 & --- & --- & 21\,843\,259.36 & --- & 21\,850\,946.24 \\
    充电量\textsuperscript{a} & 20\,740.67 & 89\,108.12 & 86\,965.89 & 88\,470.42 & 86\,423.35 \\
    放电量\textsuperscript{a} & 16\,799.94 & 72\,177.58 & 70\,442.37 & 71\,661.04 & 70\,002.91 \\
    紧急购电量 \calref{C-05} & --- & 172\,015.82 & 64\,522.09 & 182\,759.86 & 63\,523.11 \\
    盈余处置量 \calref{C-06} & 0.0000 & 3\,106\,390.84 & 2\,638\,963.35 & 3\,100\,463.79 & 2\,645\,427.81 \\
    计划购电费 & 35\,126.95 & 13\,637\,552.90 & 13\,216\,203.48 & 14\,397\,542.37 & 13\,979\,692.49 \\
    违约相关费用 \calref{C-09} & --- & --- & 210\,674.71 & --- & 224\,406.57 \\
    超额相关费用 \calref{C-09} & --- & --- & 262\,628.79 & --- & 270\,998.49 \\
    紧急购电费 & --- & 627\,239.24 & 242\,895.98 & 696\,713.47 & 247\,055.63 \\
    \textbf{总费用} & \textbf{35\,126.95} & \textbf{14\,264\,792.13} & \textbf{13\,932\,402.95} & \textbf{15\,094\,255.84} & \textbf{14\,722\,153.17} \\
    无储能基线 \(C^{base}\)（\calref{C-01}） & 48\,052.05（V3） & 16\,407\,319.63（V1） & 16\,407\,319.63（V1） & 17\,026\,782.86（V2） & 17\,026\,782.86（V2） \\
    相对 \(C^{base}\) 节省率（\calref{C-04}） & 26.90\% & 13.06\% & 15.08\% & 11.35\% & 13.54\% \\
    \bottomrule
  \end{tabularx}
  \normalsize
\end{table}

注 a：问题 1 为典型日全天值；问题 2--4-3 为题面指定四个日期（2025-03-20、06-21、
09-23、12-21）的合计值，逐日值见表~\ref{tab:q2-t2}、表~\ref{tab:key4}；334 天全窗口合计未纳入台账，
可由 \texttt{result*.xlsx} 的「充放电量」工作表逐日求和复算。

表中问题 4 的窗口级数值按 334 天求和：问题 4-3 的全天调整购电量、问题 4-2 与 4-3 的
窗口盈余处置量（\calref{C-06}），可由 \texttt{result4-2.xlsx}、\texttt{result4-3.xlsx} 与附件 2
按 \calref{C-02}、\calref{C-06} 的定义复算后补入。

表~\ref{tab:summary} 的验收：四项费用之和与总费用偏差为 0（\calref{C-03}），每列均可追溯到对应的
\texttt{result*.xlsx}。紧急购电量 \calref{C-05} 分别在问题 2、3、4-2 的 296、302、294 天触发；
问题 2 的 \(172\,015.82\)~kWh 占该窗口计划购电量的 \(0.77\%\)，对应紧急购电费占总费用 \(4.40\%\)。
固定价口径下问题 2 日均节省 \(6\,414.75\)~元（问题 1 典型日 \(12\,925.10\)~元），差额主要来自
预报误差的紧急购电费（日均 \(1\,877.96\)~元），问题 3 四次滚动后升至 \(7\,409.93\)~元；
波动价口径下同一套模型的费用整体上移，但"四次滚动优于一次计划"的方向不变。\evid{C-Q1-03}\evid{C-Q3-01}

\subsection{成本分解与基线对照}

图~\ref{fig:costwaterfall} 给出问题 1 的日费用分解：底柱为无储能基线（\calref{C-01} 变体 V3）、
中间柱为储能削峰填谷收益、顶柱为最优购电费，三者同口径、同单位（元/日）。

\begin{figure}[htbp]
  \centering
  \includefig[width=0.80\linewidth]{fig05_cost_waterfall.pdf}
  \caption{问题 1 日费用分解瀑布图（对象：微网典型日购电费用；范围：附件 1 全年平均日、
           单日 144 个 10~min 槽；基线：无储能基线 \(C^{base}=48\,052.05\)~元/日（\calref{C-01} 变体 V3）；
           中间柱：储能削峰填谷收益 \(12\,925.10\)~元/日；终点：最优购电费 \(35\,126.95\)~元/日；
           单位：元/日）}
  \label{fig:costwaterfall}
\end{figure}

334 天窗口上的同一套分解由表~\ref{tab:summary} 承载：问题 2 的计划购电费与紧急购电费分别占总费用的
\(95.60\%\) 与 \(4.40\%\)——储能的价值在压低计划购电费，紧急购电费是预报误差的代价，两者不可相抵
后只报净数。四次滚动（\calref{C-09}）把四项占比改为 \(94.86\%\)、\(1.51\%\)、\(1.89\%\)、\(1.74\%\)；
波动价口径下问题 4-2、4-3 相对基线分别节省 \(1\,932\,527.02\)~元与 \(2\,304\,629.69\)~元，
两套基线不得混用。

\textbf{策略机制。}问题 1 的最优解（图~\ref{fig:mechanism}）中，充电量 \(20\,740.67\)~kWh 落于低谷
\(0{:}00\)--\(4{:}00\)（\(4\,500.00\)~kWh）、\(20{:}00\)--\(24{:}00\)（\(5\,333.33\)~kWh）与午间
\(12{:}00\)--\(16{:}00\)（\(5\,286.04\)~kWh），放电量 \(16\,799.94\)~kWh 集中于 \(4{:}00\)--\(8{:}00\)
（\(6\,365.84\)~kWh）与 \(16{:}00\)--\(20{:}00\)（\(5\,780.13\)~kWh）；两段低谷占全日充电量的
\(47.4\%\)、午间盈余段仅占 \(25.5\%\)，放电在两段高峰合计占 \(72.3\%\)——故主机制是
\textbf{面向电价的跨时段搬运}，而非"把午间盈余存到傍晚"。\evid{C-Q1-02}

对偶量给出与求解器无关的策略形式（图~\ref{fig:dualprice}）：最优性条件为
"\(\alpha_t\le\eta_c\nu_t\) 时充电、\(\alpha_t\ge\nu_t/\eta_d\) 时放电"（\(\nu_t,\alpha_t\) 分别为储电水价值
与电量平衡影子价格），实测 \(\nu_t\in[0.4711,1.1305]\)、\(\alpha_t\in[0.3713,1.2561]\)~元/kWh，
\(E_t>0\) 时 \(\alpha_t=p_t\)（偏差 \(0.000\)）；144 个槽（充电 45、放电 37）的阈值判定
\textbf{违反次数为 0}，互补松弛最大残差 \(0.000\)，死区比 \(1/(\eta_c\eta_d)=1.2346\) 与实测充放比一致。
\evid{C-Q1-02}

问题 1 的储电量日内触及下限 \(1200.00\)~kWh 与上限 \(10\,800.00\)~kWh，\(0{:}00\) 与 \(24{:}00\)
均为 \(6000.00\)~kWh。以典型日为基准算例（容量 \(12\,000\)~kWh、功率 \(5000\)~kW、费用 \(35\,126.95\)~元），
仅解除容量上限费用降至 \(27\,351.55\)~元（再省 \(22.14\%\)），仅解除充放功率上限降至 \(34\,544.43\)~元
（再省 \(1.66\%\)）：\textbf{容量摆幅是主要瓶颈，充放电功率在部分时段同时构成约束，两者量级不同}。

问题 2 在 \(2025\text{-}03\text{-}20\) 的解复现同一结构：当日充电 \(24\,072.26\)~kWh、放电
\(19\,498.53\)~kWh，比值同为 \(1.2346\)，逐块充放电量与储电量见表~\ref{tab:q2-t2}（其中
\(8{:}00\)--\(16{:}00\) 两段充电合计 \(13\,405.59\)~kWh）；被 \(5\) 倍紧急购电覆盖的缺口集中在上午
（表~\ref{tab:emergency} 的 \(6{:}00\)--\(12{:}00\) 段），四种方案下规律一致，故安全边际与预报纠偏
应加在上午时段，而不是全天均匀加码。

\subsection{题面指定时段与日期的结果表}

按题面要求，问题 2、3、4 需对 \(2025\text{-}03\text{-}20\)、\(06\text{-}21\)、\(09\text{-}23\)、\(12\text{-}21\)
四个日期各给出表 1 格式（指定 6 个 10~min 时段购电量、全天购电量与全天购电费）、表 2 格式
（6 个 4~h 块充放电量与 \(0{:}00\)、\(24{:}00\) 储电量）与表 3 格式（紧急购电时间段与购电量）。
本文以\textbf{问题 2 为主模型}，其表 1、表 2 格式按四个日期完整给出（表~\ref{tab:q2-t1}、
表~\ref{tab:q2-t2}）；问题 3 与 4-2、4-3 不逐日铺开，由表~\ref{tab:key4} 汇总承载（逐槽数值见
\texttt{result3.xlsx}、\texttt{result4-2.xlsx}、\texttt{result4-3.xlsx}，对应关系见附录 C），
紧急购电结果按题面表 3 格式合并于表~\ref{tab:emergency}。

\begin{table}[htbp]
  \centering
  \caption{问题 2 在四个指定日期的购电量（表 1 格式；单位：购电量 kWh、购电费 元；
           时段为 10~min，按附录 5 模板的位置读法，\(10{:}00\text{-}10{:}10\) 对应附件标签 \(10{:}10\)；
           全天购电量与全天购电费为对应日期 144 个槽的合计）}
  \label{tab:q2-t1}
  \footnotesize
  \setlength{\tabcolsep}{5pt}
  \begin{tabular}{lcccc}
    \toprule
    时间段 & 2025-03-20 & 2025-06-21 & 2025-09-23 & 2025-12-21 \\
    \midrule
    \(10{:}00\)--\(10{:}10\) & 0.0000 & 0.0000 & 0.0000 & 0.0000 \\
    \(12{:}00\)--\(12{:}10\) & 804.3029 & 0.0000 & 536.2415 & 997.0335 \\
    \(14{:}00\)--\(14{:}10\) & 0.0000 & 0.0000 & 0.0000 & 0.0000 \\
    \(16{:}00\)--\(16{:}10\) & 817.3061 & 255.2254 & 902.7969 & 1\,001.2240 \\
    \(18{:}00\)--\(18{:}10\) & 697.9973 & 0.0000 & 764.3298 & 715.4144 \\
    \(20{:}00\)--\(20{:}10\) & 0.0000 & 0.0000 & 0.0000 & 0.0000 \\
    \midrule
    全天购电量（kWh） & 74\,178.1298 & 35\,701.5391 & 74\,744.7963 & 99\,692.6601 \\
    全天购电费（元） & 45\,229.90 & 20\,324.26 & 47\,063.39 & 64\,376.98 \\
    \bottomrule
  \end{tabular}
\end{table}

\begin{table}[htbp]
  \centering
  \caption{问题 2 在四个指定日期的储能充放电量与储电量（表 2 格式；单位：kWh；
           6 个 4~h 块各含 24 个 10~min 槽；\(0{:}00\) 与 \(24{:}00\) 储电量对应题面表 2 的"时刻/储电量"列，
           四个日期均为 \(6000.00\)~kWh）}
  \label{tab:q2-t2}
  \footnotesize
  \setlength{\tabcolsep}{4pt}
  \begin{tabular}{l*{4}{cc}}
    \toprule
    \multirow{2}{*}{时间段} & \multicolumn{2}{c}{2025-03-20} & \multicolumn{2}{c}{2025-06-21}
                            & \multicolumn{2}{c}{2025-09-23} & \multicolumn{2}{c}{2025-12-21} \\
    \cmidrule(lr){2-3}\cmidrule(lr){4-5}\cmidrule(lr){6-7}\cmidrule(lr){8-9}
     & 充电量 & 放电量 & 充电量 & 放电量 & 充电量 & 放电量 & 充电量 & 放电量 \\
    \midrule
    \(0{:}00\)--\(4{:}00\)   & 4\,500.0000 & 0.0000 & 0.0000 & 0.0000 & 4\,500.0000 & 0.0000 & 4\,500.0000 & 0.0000 \\
    \(4{:}00\)--\(8{:}00\)   & 833.3333 & 6\,224.1108 & 140.7851 & 4\,303.4721 & 833.3333 & 5\,547.8264 & 1\,666.6667 & 1\,508.3333 \\
    \(8{:}00\)--\(12{:}00\)  & 6\,500.0000 & 2\,415.8892 & 5\,207.8672 & 130.5638 & 5\,833.3333 & 3\,020.9227 & 5\,833.3333 & 7\,806.6667 \\
    \(12{:}00\)--\(16{:}00\) & 6\,905.5931 & 2\,218.5304 & 5\,458.7995 & 0.0000 & 6\,728.4081 & 1\,606.2615 & 8\,333.3333 & 2\,835.0000 \\
    \(16{:}00\)--\(20{:}00\) & 0.0000 & 5\,709.7835 & 0.0000 & 6\,154.1006 & 0.0000 & 5\,602.9091 & 0.0000 & 5\,797.4742 \\
    \(20{:}00\)--\(24{:}00\) & 5\,333.3333 & 2\,930.2165 & 5\,333.3333 & 2\,485.8994 & 5\,333.3333 & 3\,037.0909 & 5\,333.3333 & 2\,842.5258 \\
    \midrule
    \(0{:}00\) 储电量 & \multicolumn{2}{c}{6\,000.00} & \multicolumn{2}{c}{6\,000.00}
                      & \multicolumn{2}{c}{6\,000.00} & \multicolumn{2}{c}{6\,000.00} \\
    \(24{:}00\) 储电量 & \multicolumn{2}{c}{6\,000.00} & \multicolumn{2}{c}{6\,000.00}
                       & \multicolumn{2}{c}{6\,000.00} & \multicolumn{2}{c}{6\,000.00} \\
    \bottomrule
  \end{tabular}
\end{table}

\begin{table}[htbp]
  \centering
  \caption{四个指定日期的紧急购电时间段与购电量（表 3 格式；单位：kWh；
           行 = 日期 \(\times\) 方案，括号内为该时段购电量，末列为该日合计；
           时段按连续触发的 10~min 槽合并、左闭右开；
           问题 2、4-2 无调整机制，问题 4-2 在 \(2025\text{-}12\text{-}21\) 未触发紧急购电）}
  \label{tab:emergency}
  \footnotesize
  \setlength{\tabcolsep}{4pt}
  \begin{tabular}{C{1.35cm} C{1.85cm} L{8.10cm} C{1.85cm}}
    \toprule
    方案 & 日期 & 紧急购电时间段（括号内购电量，kWh） & 合计（kWh） \\
    \midrule
    问题 2 & 2025-03-20 & \(6{:}00\)--\(8{:}00\)（228.99）；\(8{:}10\)--\(8{:}30\)（49.16）；\(8{:}50\)--\(10{:}30\)（531.27）；\(10{:}40\)--\(11{:}40\)（216.46） & 1\,025.88 \\
    问题 2 & 2025-06-21 & \(4{:}00\)--\(4{:}50\)（5.24）；\(14{:}00\)--\(14{:}30\)（181.84） & 187.08 \\
    问题 2 & 2025-09-23 & \(6{:}50\)--\(7{:}40\)（132.59）；\(11{:}00\)--\(11{:}40\)（60.77） & 193.37 \\
    问题 2 & 2025-12-21 & \(11{:}00\)--\(11{:}30\)（22.53） & 22.53 \\
    \midrule
    问题 3 & 2025-03-20 & \(6{:}00\)--\(7{:}00\)（138.11）；\(7{:}40\)--\(9{:}20\)（206.13）；\(9{:}30\)--\(10{:}00\)（88.68）；\(10{:}50\)--\(12{:}00\)（344.46）；\(12{:}30\)--\(12{:}50\)（102.31）；\(13{:}20\)--\(13{:}50\)（94.60） & 974.29 \\
    问题 3 & 2025-06-21 & \(4{:}00\)--\(4{:}50\)（5.24） & 5.24 \\
    问题 3 & 2025-09-23 & \(7{:}00\)--\(7{:}40\)（52.76）；\(8{:}30\)--\(8{:}40\)（4.26）；\(11{:}00\)--\(11{:}10\)（40.45） & 97.47 \\
    问题 3 & 2025-12-21 & \(9{:}50\)--\(10{:}00\)（2.12）；\(11{:}00\)--\(12{:}00\)（358.56） & 360.69 \\
    \midrule
    问题 4-2 & 2025-03-20 & \(6{:}00\)--\(8{:}00\)（225.30）；\(8{:}10\)--\(8{:}40\)（55.07）；\(8{:}50\)--\(10{:}30\)（523.99）；\(10{:}40\)--\(11{:}20\)（109.82） & 914.18 \\
    问题 4-2 & 2025-06-21 & \(4{:}00\)--\(4{:}50\)（5.24）；\(5{:}00\)--\(5{:}30\)（82.79）；\(14{:}00\)--\(14{:}30\)（164.79） & 252.82 \\
    问题 4-2 & 2025-09-23 & \(6{:}20\)--\(6{:}30\)（1.30）；\(6{:}50\)--\(7{:}40\)（125.27）；\(11{:}00\)--\(11{:}10\)（23.07） & 149.64 \\
    问题 4-2 & 2025-12-21 & 无 & 0.00 \\
    \midrule
    问题 4-3 & 2025-03-20 & \(6{:}00\)--\(7{:}00\)（125.36）；\(7{:}40\)--\(9{:}20\)（181.60）；\(9{:}30\)--\(10{:}00\)（87.91）；\(10{:}50\)--\(12{:}00\)（325.44）；\(12{:}30\)--\(12{:}50\)（86.88）；\(13{:}20\)--\(13{:}40\)（87.36） & 894.55 \\
    问题 4-3 & 2025-06-21 & \(4{:}00\)--\(4{:}50\)（5.24） & 5.24 \\
    问题 4-3 & 2025-09-23 & \(7{:}00\)--\(7{:}40\)（61.40）；\(11{:}00\)--\(11{:}10\)（30.51） & 91.92 \\
    问题 4-3 & 2025-12-21 & \(11{:}00\)--\(12{:}00\)（345.33） & 345.33 \\
    \bottomrule
  \end{tabular}
\end{table}

\begin{table}[htbp]
  \centering
  \caption{问题 3、4-2、4-3 在四个指定日期的汇总结果（单位：电量 kWh、费用 元；
           全部为当日 144 个槽的合计；"全天购电费"为 \calref{C-09} 偏差量定价口径下的 \(C^{\text{plan}}\)；
           该日总费用 \(=\) 全天购电费 \(+\) 违约与超额相关费用 \(+\) 紧急购电费）}
  \label{tab:key4}
  \footnotesize
  \setlength{\tabcolsep}{3.5pt}
  \begin{tabular}{C{1.35cm} C{1.75cm} C{2.10cm} C{2.10cm} C{1.85cm} C{1.75cm} C{1.75cm}}
    \toprule
    方案 & 日期 & 全天计划购电量 & 全天调整购电量 & 全天购电费 & 紧急购电量 & 该日总费用 \\
    \midrule
    问题 3 & 2025-03-20 & 74\,178.13 & 69\,212.46 & 41\,926.08 & 974.29 & 47\,482.28 \\
    问题 3 & 2025-06-21 & 35\,701.54 & 35\,432.45 & 20\,103.36 & 5.24 & 20\,260.50 \\
    问题 3 & 2025-09-23 & 74\,744.80 & 72\,213.13 & 45\,415.68 & 97.47 & 46\,558.94 \\
    问题 3 & 2025-12-21 & 99\,692.66 & 98\,984.11 & 63\,729.13 & 360.69 & 65\,154.58 \\
    \midrule
    问题 4-2 & 2025-03-20 & 74\,196.33 & --- & 46\,622.31 & 914.18 & 50\,923.43 \\
    问题 4-2 & 2025-06-21 & 35\,543.98 & --- & 17\,764.30 & 252.82 & 18\,587.09 \\
    问题 4-2 & 2025-09-23 & 74\,693.23 & --- & 49\,288.50 & 149.64 & 49\,839.31 \\
    问题 4-2 & 2025-12-21 & 99\,538.07 & --- & 76\,013.33 & 0.00 & 76\,013.33 \\
    \midrule
    问题 4-3 & 2025-03-20 & 74\,321.58 & 69\,229.87 & 43\,088.31 & 894.55 & 48\,781.27 \\
    问题 4-3 & 2025-06-21 & 35\,725.56 & 35\,440.19 & 17\,651.83 & 5.24 & 17\,769.08 \\
    问题 4-3 & 2025-09-23 & 74\,869.82 & 72\,305.05 & 47\,748.10 & 91.92 & 48\,876.69 \\
    问题 4-3 & 2025-12-21 & 99\,720.87 & 98\,953.92 & 75\,242.01 & 345.33 & 77\,242.42 \\
    \bottomrule
  \end{tabular}
\end{table}

四个日期中，\(2025\text{-}03\text{-}20\) 与 \(09\text{-}23\) 的负载形状接近而紧急购电量相差
\(1\,025.88\) 对 \(193.37\)~kWh，差异来自当日光伏预报偏差；\(2025\text{-}12\text{-}21\) 购电量最大
（\(99\,692.66\)~kWh）而紧急购电量最小，说明该日预报偏差最小、计划量本身已足够。四次滚动把
\(2025\text{-}03\text{-}20\) 的紧急购电量由 \(1\,025.88\)~kWh 降到 \(974.29\)~kWh，代价是
\(1\,949.15\)~元的违约与超额相关费用（\(1\,651.91+297.24\)），该日总费用净降 \(2\,171.74\)~元。
全窗口上四次滚动的净收益为 \(332\,389.18\)~元，而继续加密预报的剩余空间只有 \(58\,062.56\)~元
（第九章表~\ref{tab:ablation}）：滚动决策的收益集中在前几次预报，边际收益快速衰减。

回读校验：结果文件的工作表名、列序与模板一致；\texttt{result2.xlsx} 的「计划购电量」为
335 行 \(\times\) 147 列（334 个日期行），「充放电量」按 \(334\times6=2004\) 行重建，
「紧急购电量」按事件逐行生成（问题 2、3、4-2、4-3 分别为 1156、1190、1193、1187 行）；
全天汇总列与逐时段求和的最大偏差 \(0.0012\)~kWh，全天购电费最大相对偏差 \(4.32\times10^{-7}\)
（第九章审计 A9、A10）。
